%% =====================================================================
%%  preamble.tex -- zajednička pravila formatiranja za cijelu dokumentaciju
%%  Učitava se iz main.tex i iz AI izvještaja (%% =====================================================================
%%  preamble.tex -- zajednička pravila formatiranja za cijelu dokumentaciju
%%  Učitava se iz main.tex i iz AI izvještaja (%% =====================================================================
%%  preamble.tex -- zajednička pravila formatiranja za cijelu dokumentaciju
%%  Učitava se iz main.tex i iz AI izvještaja (%% =====================================================================
%%  preamble.tex -- zajednička pravila formatiranja za cijelu dokumentaciju
%%  Učitava se iz main.tex i iz AI izvještaja (\input{preamble}),
%%  pa su oba dokumenta uvijek jednako oblikovana.
%%
%%  PRAVILA ZA SVE AUTORE:
%%   * Svako poglavlje je zasebna .tex datoteka (bez \documentclass i
%%     \begin{document}) koja počinje s \section{...}.
%%   * Slike i tabele se numerišu po poglavljima (Slika 2.1, Tabela 2.3).
%%   * Potpis (\caption) ide ISPOD slike i ISPOD tabele; uvijek dodati
%%     \label (slike: fig:..., tabele: tab:...) i pozivati se s \ref{...}.
%%   * Slike i tabele ubacivati u okruženju [H] (paket float).
%%   * Literatura: \cite{kljuc}; nove stavke dodati na kraju main.tex.
%%   * Decimalni zarez u tekstu i tabelama (0,85), nazivi metrika ostaju
%%     kao u literaturi (mAP@0.5).
%%   * Crvene oznake \dopuniti{...} i \provjeriti{...} su podsjetnici;
%%     prije predaje pretražiti "dopuniti" i "provjeriti".
%% =====================================================================

\usepackage[utf8]{inputenc}
\usepackage[T1]{fontenc}
\usepackage{lmodern}
% bosnian zahtijeva noviju TeX Live distribuciju; ako nije dostupan, koristi se croatian
\IfFileExists{bosnian.ldf}{%
  \usepackage[bosnian]{babel}%
  \addto\captionsbosnian{\renewcommand{\tablename}{Tabela}\renewcommand{\refname}{Literatura}\renewcommand{\contentsname}{Sadržaj}}%
}{%
  \usepackage[croatian]{babel}%
  \addto\captionscroatian{\renewcommand{\tablename}{Tabela}\renewcommand{\refname}{Literatura}\renewcommand{\contentsname}{Sadržaj}}%
}

\usepackage[a4paper,margin=2.5cm]{geometry}
\usepackage{microtype}
\usepackage{amsmath,amssymb}
\usepackage{booktabs,array,tabularx}
\usepackage{float}
\usepackage{caption}
\usepackage{enumitem}
\usepackage{xcolor}
\usepackage{tcolorbox}
\usepackage{tikz}
\usetikzlibrary{positioning,arrows.meta,shapes.geometric,calc}
\usepackage{pgfplots}
\pgfplotsset{compat=1.17}
\usepackage{fancyhdr}
\usepackage{titlesec}
\usepackage[hidelinks]{hyperref}

% ---------- numeracija slika, tabela i jednačina po poglavljima ----------
\counterwithin{figure}{section}
\counterwithin{table}{section}
\counterwithin{equation}{section}

% ---------- boje i oznake ----------
\definecolor{etfblue}{RGB}{0,51,102}
\newcommand{\dopuniti}[1]{\textcolor{red}{\textbf{[#1]}}}
\newcommand{\provjeriti}[1]{\par\smallskip{\small\textcolor{red}{\textit{Napomena (ukloniti nakon provjere): #1}}}\par}

% ---------- podaci o projektu i timu (mijenjati samo ovdje) ----------
\newcommand{\NazivTima}{Vatrogasci}
\newcommand{\NazivProjekta}{Detekcija vatre i dima na slikama}
\newcommand{\PodnaslovProjekta}{Razlikovanje prisutnosti dima i vatre na skupu podataka D-Fire}
\newcommand{\AutorA}{Danijal Alibegović}
\newcommand{\AutorB}{Aldin Velić}
\newcommand{\AutorC}{Zijad Doglod}
\newcommand{\AutorD}{Dino Burić}
\newcommand{\AutorE}{Haris Macić}
\newcommand{\AutorF}{Adnan Šemić}

% ---------- osnovno oblikovanje ----------
\setlength{\parindent}{0pt}
\setlength{\parskip}{5pt plus 1pt minus 1pt}
\renewcommand{\arraystretch}{1.2}
\emergencystretch=1.5em
\captionsetup{font=small,labelfont=bf,labelsep=period,skip=6pt}
\setlist{topsep=3pt,itemsep=2pt,parsep=0pt,leftmargin=*}
\setcounter{tocdepth}{2}

\titleformat{\section}{\Large\bfseries\color{etfblue}}{\thesection.}{0.6em}{}[\color{etfblue}\titlerule]
\titleformat{name=\section,numberless}{\Large\bfseries\color{etfblue}}{}{0pt}{}[\color{etfblue}\titlerule]
\titleformat{\subsection}{\large\bfseries\color{etfblue}}{\thesubsection}{0.6em}{}
\titleformat{name=\subsection,numberless}{\large\bfseries\color{etfblue}}{}{0pt}{}
\titleformat{\subsubsection}{\normalsize\bfseries}{\thesubsubsection}{0.6em}{}
\titlespacing*{\section}{0pt}{16pt}{8pt}
\titlespacing*{\subsection}{0pt}{12pt}{4pt}
\titlespacing*{\subsubsection}{0pt}{8pt}{2pt}

\pagestyle{fancy}
\fancyhf{}
\fancyhead[L]{\small ETF Sarajevo $\bullet$ POOS $\bullet$ Tim \NazivTima}
\fancyhead[R]{\small Projektni zadatak}
\fancyfoot[C]{\small\thepage}
\renewcommand{\headrulewidth}{0.4pt}

% ---------- novo poglavlje uvijek počinje na novoj stranici ----------
\newcommand{\poglavlje}[1]{\clearpage\input{#1}}
\newcommand{\paperbox}[5]{%
  \begin{tcolorbox}[colback=black!4,colframe=black!45,boxrule=0.4pt,arc=2pt,
      left=6pt,right=6pt,top=4pt,bottom=4pt]
  \small
  \begin{tabularx}{\linewidth}{@{}>{\bfseries}p{3.4cm}X@{}}
    Autor izvještaja: & #1 \\
    Analizirani rad: & #2 \\
    Izvor: & #3 \\
    Skup podataka u radu: & #4 \\
    Vrsta rada: & #5 \\
  \end{tabularx}
  \end{tcolorbox}\par\smallskip}),
%%  pa su oba dokumenta uvijek jednako oblikovana.
%%
%%  PRAVILA ZA SVE AUTORE:
%%   * Svako poglavlje je zasebna .tex datoteka (bez \documentclass i
%%     \begin{document}) koja počinje s \section{...}.
%%   * Slike i tabele se numerišu po poglavljima (Slika 2.1, Tabela 2.3).
%%   * Potpis (\caption) ide ISPOD slike i ISPOD tabele; uvijek dodati
%%     \label (slike: fig:..., tabele: tab:...) i pozivati se s \ref{...}.
%%   * Slike i tabele ubacivati u okruženju [H] (paket float).
%%   * Literatura: \cite{kljuc}; nove stavke dodati na kraju main.tex.
%%   * Decimalni zarez u tekstu i tabelama (0,85), nazivi metrika ostaju
%%     kao u literaturi (mAP@0.5).
%%   * Crvene oznake \dopuniti{...} i \provjeriti{...} su podsjetnici;
%%     prije predaje pretražiti "dopuniti" i "provjeriti".
%% =====================================================================

\usepackage[utf8]{inputenc}
\usepackage[T1]{fontenc}
\usepackage{lmodern}
% bosnian zahtijeva noviju TeX Live distribuciju; ako nije dostupan, koristi se croatian
\IfFileExists{bosnian.ldf}{%
  \usepackage[bosnian]{babel}%
  \addto\captionsbosnian{\renewcommand{\tablename}{Tabela}\renewcommand{\refname}{Literatura}\renewcommand{\contentsname}{Sadržaj}}%
}{%
  \usepackage[croatian]{babel}%
  \addto\captionscroatian{\renewcommand{\tablename}{Tabela}\renewcommand{\refname}{Literatura}\renewcommand{\contentsname}{Sadržaj}}%
}

\usepackage[a4paper,margin=2.5cm]{geometry}
\usepackage{microtype}
\usepackage{amsmath,amssymb}
\usepackage{booktabs,array,tabularx}
\usepackage{float}
\usepackage{caption}
\usepackage{enumitem}
\usepackage{xcolor}
\usepackage{tcolorbox}
\usepackage{tikz}
\usetikzlibrary{positioning,arrows.meta,shapes.geometric,calc}
\usepackage{pgfplots}
\pgfplotsset{compat=1.17}
\usepackage{fancyhdr}
\usepackage{titlesec}
\usepackage[hidelinks]{hyperref}

% ---------- numeracija slika, tabela i jednačina po poglavljima ----------
\counterwithin{figure}{section}
\counterwithin{table}{section}
\counterwithin{equation}{section}

% ---------- boje i oznake ----------
\definecolor{etfblue}{RGB}{0,51,102}
\newcommand{\dopuniti}[1]{\textcolor{red}{\textbf{[#1]}}}
\newcommand{\provjeriti}[1]{\par\smallskip{\small\textcolor{red}{\textit{Napomena (ukloniti nakon provjere): #1}}}\par}

% ---------- podaci o projektu i timu (mijenjati samo ovdje) ----------
\newcommand{\NazivTima}{Vatrogasci}
\newcommand{\NazivProjekta}{Detekcija vatre i dima na slikama}
\newcommand{\PodnaslovProjekta}{Razlikovanje prisutnosti dima i vatre na skupu podataka D-Fire}
\newcommand{\AutorA}{Danijal Alibegović}
\newcommand{\AutorB}{Aldin Velić}
\newcommand{\AutorC}{Zijad Doglod}
\newcommand{\AutorD}{Dino Burić}
\newcommand{\AutorE}{Haris Macić}
\newcommand{\AutorF}{Adnan Šemić}

% ---------- osnovno oblikovanje ----------
\setlength{\parindent}{0pt}
\setlength{\parskip}{5pt plus 1pt minus 1pt}
\renewcommand{\arraystretch}{1.2}
\emergencystretch=1.5em
\captionsetup{font=small,labelfont=bf,labelsep=period,skip=6pt}
\setlist{topsep=3pt,itemsep=2pt,parsep=0pt,leftmargin=*}
\setcounter{tocdepth}{2}

\titleformat{\section}{\Large\bfseries\color{etfblue}}{\thesection.}{0.6em}{}[\color{etfblue}\titlerule]
\titleformat{name=\section,numberless}{\Large\bfseries\color{etfblue}}{}{0pt}{}[\color{etfblue}\titlerule]
\titleformat{\subsection}{\large\bfseries\color{etfblue}}{\thesubsection}{0.6em}{}
\titleformat{name=\subsection,numberless}{\large\bfseries\color{etfblue}}{}{0pt}{}
\titleformat{\subsubsection}{\normalsize\bfseries}{\thesubsubsection}{0.6em}{}
\titlespacing*{\section}{0pt}{16pt}{8pt}
\titlespacing*{\subsection}{0pt}{12pt}{4pt}
\titlespacing*{\subsubsection}{0pt}{8pt}{2pt}

\pagestyle{fancy}
\fancyhf{}
\fancyhead[L]{\small ETF Sarajevo $\bullet$ POOS $\bullet$ Tim \NazivTima}
\fancyhead[R]{\small Projektni zadatak}
\fancyfoot[C]{\small\thepage}
\renewcommand{\headrulewidth}{0.4pt}

% ---------- novo poglavlje uvijek počinje na novoj stranici ----------
\newcommand{\poglavlje}[1]{\clearpage\input{#1}}
\newcommand{\paperbox}[5]{%
  \begin{tcolorbox}[colback=black!4,colframe=black!45,boxrule=0.4pt,arc=2pt,
      left=6pt,right=6pt,top=4pt,bottom=4pt]
  \small
  \begin{tabularx}{\linewidth}{@{}>{\bfseries}p{3.4cm}X@{}}
    Autor izvještaja: & #1 \\
    Analizirani rad: & #2 \\
    Izvor: & #3 \\
    Skup podataka u radu: & #4 \\
    Vrsta rada: & #5 \\
  \end{tabularx}
  \end{tcolorbox}\par\smallskip}),
%%  pa su oba dokumenta uvijek jednako oblikovana.
%%
%%  PRAVILA ZA SVE AUTORE:
%%   * Svako poglavlje je zasebna .tex datoteka (bez \documentclass i
%%     \begin{document}) koja počinje s \section{...}.
%%   * Slike i tabele se numerišu po poglavljima (Slika 2.1, Tabela 2.3).
%%   * Potpis (\caption) ide ISPOD slike i ISPOD tabele; uvijek dodati
%%     \label (slike: fig:..., tabele: tab:...) i pozivati se s \ref{...}.
%%   * Slike i tabele ubacivati u okruženju [H] (paket float).
%%   * Literatura: \cite{kljuc}; nove stavke dodati na kraju main.tex.
%%   * Decimalni zarez u tekstu i tabelama (0,85), nazivi metrika ostaju
%%     kao u literaturi (mAP@0.5).
%%   * Crvene oznake \dopuniti{...} i \provjeriti{...} su podsjetnici;
%%     prije predaje pretražiti "dopuniti" i "provjeriti".
%% =====================================================================

\usepackage[utf8]{inputenc}
\usepackage[T1]{fontenc}
\usepackage{lmodern}
% bosnian zahtijeva noviju TeX Live distribuciju; ako nije dostupan, koristi se croatian
\IfFileExists{bosnian.ldf}{%
  \usepackage[bosnian]{babel}%
  \addto\captionsbosnian{\renewcommand{\tablename}{Tabela}\renewcommand{\refname}{Literatura}\renewcommand{\contentsname}{Sadržaj}}%
}{%
  \usepackage[croatian]{babel}%
  \addto\captionscroatian{\renewcommand{\tablename}{Tabela}\renewcommand{\refname}{Literatura}\renewcommand{\contentsname}{Sadržaj}}%
}

\usepackage[a4paper,margin=2.5cm]{geometry}
\usepackage{microtype}
\usepackage{amsmath,amssymb}
\usepackage{booktabs,array,tabularx}
\usepackage{float}
\usepackage{caption}
\usepackage{enumitem}
\usepackage{xcolor}
\usepackage{tcolorbox}
\usepackage{tikz}
\usetikzlibrary{positioning,arrows.meta,shapes.geometric,calc}
\usepackage{pgfplots}
\pgfplotsset{compat=1.17}
\usepackage{fancyhdr}
\usepackage{titlesec}
\usepackage[hidelinks]{hyperref}

% ---------- numeracija slika, tabela i jednačina po poglavljima ----------
\counterwithin{figure}{section}
\counterwithin{table}{section}
\counterwithin{equation}{section}

% ---------- boje i oznake ----------
\definecolor{etfblue}{RGB}{0,51,102}
\newcommand{\dopuniti}[1]{\textcolor{red}{\textbf{[#1]}}}
\newcommand{\provjeriti}[1]{\par\smallskip{\small\textcolor{red}{\textit{Napomena (ukloniti nakon provjere): #1}}}\par}

% ---------- podaci o projektu i timu (mijenjati samo ovdje) ----------
\newcommand{\NazivTima}{Vatrogasci}
\newcommand{\NazivProjekta}{Detekcija vatre i dima na slikama}
\newcommand{\PodnaslovProjekta}{Razlikovanje prisutnosti dima i vatre na skupu podataka D-Fire}
\newcommand{\AutorA}{Danijal Alibegović}
\newcommand{\AutorB}{Aldin Velić}
\newcommand{\AutorC}{Zijad Doglod}
\newcommand{\AutorD}{Dino Burić}
\newcommand{\AutorE}{Haris Macić}
\newcommand{\AutorF}{Adnan Šemić}

% ---------- osnovno oblikovanje ----------
\setlength{\parindent}{0pt}
\setlength{\parskip}{5pt plus 1pt minus 1pt}
\renewcommand{\arraystretch}{1.2}
\emergencystretch=1.5em
\captionsetup{font=small,labelfont=bf,labelsep=period,skip=6pt}
\setlist{topsep=3pt,itemsep=2pt,parsep=0pt,leftmargin=*}
\setcounter{tocdepth}{2}

\titleformat{\section}{\Large\bfseries\color{etfblue}}{\thesection.}{0.6em}{}[\color{etfblue}\titlerule]
\titleformat{name=\section,numberless}{\Large\bfseries\color{etfblue}}{}{0pt}{}[\color{etfblue}\titlerule]
\titleformat{\subsection}{\large\bfseries\color{etfblue}}{\thesubsection}{0.6em}{}
\titleformat{name=\subsection,numberless}{\large\bfseries\color{etfblue}}{}{0pt}{}
\titleformat{\subsubsection}{\normalsize\bfseries}{\thesubsubsection}{0.6em}{}
\titlespacing*{\section}{0pt}{16pt}{8pt}
\titlespacing*{\subsection}{0pt}{12pt}{4pt}
\titlespacing*{\subsubsection}{0pt}{8pt}{2pt}

\pagestyle{fancy}
\fancyhf{}
\fancyhead[L]{\small ETF Sarajevo $\bullet$ POOS $\bullet$ Tim \NazivTima}
\fancyhead[R]{\small Projektni zadatak}
\fancyfoot[C]{\small\thepage}
\renewcommand{\headrulewidth}{0.4pt}

% ---------- novo poglavlje uvijek počinje na novoj stranici ----------
\newcommand{\poglavlje}[1]{\clearpage\input{#1}}
\newcommand{\paperbox}[5]{%
  \begin{tcolorbox}[colback=black!4,colframe=black!45,boxrule=0.4pt,arc=2pt,
      left=6pt,right=6pt,top=4pt,bottom=4pt]
  \small
  \begin{tabularx}{\linewidth}{@{}>{\bfseries}p{3.4cm}X@{}}
    Autor izvještaja: & #1 \\
    Analizirani rad: & #2 \\
    Izvor: & #3 \\
    Skup podataka u radu: & #4 \\
    Vrsta rada: & #5 \\
  \end{tabularx}
  \end{tcolorbox}\par\smallskip}),
%%  pa su oba dokumenta uvijek jednako oblikovana.
%%
%%  PRAVILA ZA SVE AUTORE:
%%   * Svako poglavlje je zasebna .tex datoteka (bez \documentclass i
%%     \begin{document}) koja počinje s \section{...}.
%%   * Slike i tabele se numerišu po poglavljima (Slika 2.1, Tabela 2.3).
%%   * Potpis (\caption) ide ISPOD slike i ISPOD tabele; uvijek dodati
%%     \label (slike: fig:..., tabele: tab:...) i pozivati se s \ref{...}.
%%   * Slike i tabele ubacivati u okruženju [H] (paket float).
%%   * Literatura: \cite{kljuc}; nove stavke dodati na kraju main.tex.
%%   * Decimalni zarez u tekstu i tabelama (0,85), nazivi metrika ostaju
%%     kao u literaturi (mAP@0.5).
%%   * Crvene oznake \dopuniti{...} i \provjeriti{...} su podsjetnici;
%%     prije predaje pretražiti "dopuniti" i "provjeriti".
%% =====================================================================

\usepackage[utf8]{inputenc}
\usepackage[T1]{fontenc}
\usepackage{lmodern}
% bosnian zahtijeva noviju TeX Live distribuciju; ako nije dostupan, koristi se croatian
\IfFileExists{bosnian.ldf}{%
  \usepackage[bosnian]{babel}%
  \addto\captionsbosnian{\renewcommand{\tablename}{Tabela}\renewcommand{\refname}{Literatura}\renewcommand{\contentsname}{Sadržaj}}%
}{%
  \usepackage[croatian]{babel}%
  \addto\captionscroatian{\renewcommand{\tablename}{Tabela}\renewcommand{\refname}{Literatura}\renewcommand{\contentsname}{Sadržaj}}%
}

\usepackage[a4paper,margin=2.5cm]{geometry}
\usepackage{microtype}
\usepackage{amsmath,amssymb}
\usepackage{booktabs,array,tabularx}
\usepackage{float}
\usepackage{caption}
\usepackage{enumitem}
\usepackage{xcolor}
\usepackage{tcolorbox}
\usepackage{tikz}
\usetikzlibrary{positioning,arrows.meta,shapes.geometric,calc}
\usepackage{pgfplots}
\pgfplotsset{compat=1.17}
\usepackage{fancyhdr}
\usepackage{titlesec}
\usepackage[hidelinks]{hyperref}

% ---------- numeracija slika, tabela i jednačina po poglavljima ----------
\counterwithin{figure}{section}
\counterwithin{table}{section}
\counterwithin{equation}{section}

% ---------- boje i oznake ----------
\definecolor{etfblue}{RGB}{0,51,102}
\newcommand{\dopuniti}[1]{\textcolor{red}{\textbf{[#1]}}}
\newcommand{\provjeriti}[1]{\par\smallskip{\small\textcolor{red}{\textit{Napomena (ukloniti nakon provjere): #1}}}\par}

% ---------- podaci o projektu i timu (mijenjati samo ovdje) ----------
\newcommand{\NazivTima}{Vatrogasci}
\newcommand{\NazivProjekta}{Detekcija vatre i dima na slikama}
\newcommand{\PodnaslovProjekta}{Razlikovanje prisutnosti dima i vatre na skupu podataka D-Fire}
\newcommand{\AutorA}{Danijal Alibegović}
\newcommand{\AutorB}{Aldin Velić}
\newcommand{\AutorC}{Zijad Doglod}
\newcommand{\AutorD}{Dino Burić}
\newcommand{\AutorE}{Haris Macić}
\newcommand{\AutorF}{Adnan Šemić}

% ---------- osnovno oblikovanje ----------
\setlength{\parindent}{0pt}
\setlength{\parskip}{5pt plus 1pt minus 1pt}
\renewcommand{\arraystretch}{1.2}
\emergencystretch=1.5em
\captionsetup{font=small,labelfont=bf,labelsep=period,skip=6pt}
\setlist{topsep=3pt,itemsep=2pt,parsep=0pt,leftmargin=*}
\setcounter{tocdepth}{2}

\titleformat{\section}{\Large\bfseries\color{etfblue}}{\thesection.}{0.6em}{}[\color{etfblue}\titlerule]
\titleformat{name=\section,numberless}{\Large\bfseries\color{etfblue}}{}{0pt}{}[\color{etfblue}\titlerule]
\titleformat{\subsection}{\large\bfseries\color{etfblue}}{\thesubsection}{0.6em}{}
\titleformat{name=\subsection,numberless}{\large\bfseries\color{etfblue}}{}{0pt}{}
\titleformat{\subsubsection}{\normalsize\bfseries}{\thesubsubsection}{0.6em}{}
\titlespacing*{\section}{0pt}{16pt}{8pt}
\titlespacing*{\subsection}{0pt}{12pt}{4pt}
\titlespacing*{\subsubsection}{0pt}{8pt}{2pt}

\pagestyle{fancy}
\fancyhf{}
\fancyhead[L]{\small ETF Sarajevo $\bullet$ POOS $\bullet$ Tim \NazivTima}
\fancyhead[R]{\small Projektni zadatak}
\fancyfoot[C]{\small\thepage}
\renewcommand{\headrulewidth}{0.4pt}

% ---------- novo poglavlje uvijek počinje na novoj stranici ----------
\newcommand{\poglavlje}[1]{\clearpage\input{#1}}
\newcommand{\paperbox}[5]{%
  \begin{tcolorbox}[colback=black!4,colframe=black!45,boxrule=0.4pt,arc=2pt,
      left=6pt,right=6pt,top=4pt,bottom=4pt]
  \small
  \begin{tabularx}{\linewidth}{@{}>{\bfseries}p{3.4cm}X@{}}
    Autor izvještaja: & #1 \\
    Analizirani rad: & #2 \\
    Izvor: & #3 \\
    Skup podataka u radu: & #4 \\
    Vrsta rada: & #5 \\
  \end{tabularx}
  \end{tcolorbox}\par\smallskip}